% !TEX program = xelatex
% 第二大脑系统架构总结 —— 脱敏版
\documentclass[11pt,a4paper]{ctexart}

\usepackage[margin=2.3cm,top=2.4cm,bottom=2.4cm]{geometry}
\usepackage{booktabs}
\usepackage{longtable}
\usepackage{array}
\usepackage{xcolor}
\usepackage{enumitem}
\usepackage{amssymb}
\usepackage{titlesec}
\usepackage{fancyhdr}
\usepackage{ragged2e}
\usepackage[hidelinks]{hyperref}

% ---------- 配色 ----------
\definecolor{MainNavy}{HTML}{1F3864}
\definecolor{AccentTeal}{HTML}{2E7D8F}
\definecolor{WarnRed}{HTML}{A62B2B}
\definecolor{SoftGray}{HTML}{F2F4F7}
\definecolor{LineGray}{HTML}{9AA5B1}
\definecolor{QuoteBg}{HTML}{F7F9FB}

% ---------- 标题样式 ----------
\titleformat{\section}
  {\normalfont\Large\bfseries\color{MainNavy}}{\thesection}{0.7em}{}
\titleformat{\subsection}
  {\normalfont\large\bfseries\color{AccentTeal}}{\thesubsection}{0.6em}{}
\titleformat{\subsubsection}
  {\normalfont\normalsize\bfseries\color{MainNavy}}{\thesubsubsection}{0.5em}{}
\titlespacing*{\section}{0pt}{1.6em}{0.7em}
\titlespacing*{\subsection}{0pt}{1.1em}{0.5em}

% ---------- 页眉页脚 ----------
\pagestyle{fancy}
\fancyhf{}
\fancyhead[L]{\small\color{LineGray}第二大脑系统架构总结}
\fancyhead[R]{\small\color{LineGray}脱敏版}
\fancyfoot[C]{\small\thepage}
\renewcommand{\headrulewidth}{0.4pt}
\renewcommand{\headrule}{\hbox to\headwidth{\color{LineGray}\leaders\hrule height \headrulewidth\hfill}}

% ---------- 自定义盒子 ----------
\newcommand{\keybox}[1]{%
  \begin{center}%
  \begin{minipage}{0.94\linewidth}%
  \setlength{\fboxsep}{9pt}%
  \colorbox{SoftGray}{\parbox{\dimexpr\linewidth-2\fboxsep}{#1}}%
  \end{minipage}%
  \end{center}%
}

\newcommand{\warnbox}[1]{%
  \begin{center}%
  \begin{minipage}{0.94\linewidth}%
  \setlength{\fboxsep}{9pt}%
  \colorbox{QuoteBg}{\parbox{\dimexpr\linewidth-2\fboxsep}{%
    \color{WarnRed}#1}}%
  \end{minipage}%
  \end{center}%
}

% 表格行距紧凑
\renewcommand{\arraystretch}{1.25}
\setlength{\tabcolsep}{5pt}

% 列表间距
\setlist[itemize]{leftmargin=1.5em,itemsep=2pt,topsep=3pt,parsep=0pt}
\setlist[enumerate]{leftmargin=1.7em,itemsep=2pt,topsep=3pt,parsep=0pt}

\setlength{\emergencystretch}{3em}

% ======================================================================
\begin{document}

\begin{titlepage}
\centering
\vspace*{2.6cm}
{\color{MainNavy}\rule{\linewidth}{1.6pt}}\\[0.9em]
{\Huge\bfseries\color{MainNavy} 第二大脑系统架构总结}\\[0.55em]
{\large\color{AccentTeal} 运转机制 · 核心亮点 · 真实缺陷}\\[0.9em]
{\color{MainNavy}\rule{\linewidth}{1.6pt}}\\[2.6em]

\begin{minipage}{0.84\linewidth}
\centering
\large
\begin{tabular}{@{}r@{\hspace{1.2em}}l@{}}
\textbf{文档性质} & 系统框架总结与缺陷评估 \\[0.35em]
\textbf{描述口径} & 只涉及整体架构，不涉及具体项目 \\
\textbf{数据基准} & 系统自述的最近一次健康检测快照 \\
\textbf{数据来源} & 系统自身度量产物与运行日志 \\
\textbf{脱敏状态} & 路径 / 标识 / 设备 / 个人信息均已剔除 \\
\end{tabular}
\end{minipage}

\vfill
\keybox{\small\centering
\textbf{本文所有数字均为系统自身产物的原值，未作估算或修饰。}}
\vspace{1.4cm}
\end{titlepage}

\tableofcontents
\newpage

% ======================================================================
\section{系统定位}

这是一套面向多个 AI Agent 的\textbf{共用长期记忆系统}。它的核心命题不是「给某个 AI 加一个知识库」，而是解决四个具体问题：

\begin{center}
\begin{tabular}{@{}p{0.30\linewidth}p{0.31\linewidth}p{0.31\linewidth}@{}}
\toprule
\textbf{问题} & \textbf{传统做法的失效点} & \textbf{本系统的应对} \\
\midrule
每次新会话从零开始 & 上下文窗口关闭即遗忘 & 会话启动时自动注入压缩后的记忆包 \\
多个 AI 之间互不知情 & 各自为政，同一个坑重复踩 & 单一共享仓库 + 跨 Agent 独立验证 \\
说过的偏好反复重申 & 偏好散落在历史对话里 & 偏好升格为受治理的正式条目 \\
Agent 自以为做完 & 无回执、无版本、无审计 & 写入登记 + 版本号 + 幂等日志 \\
\bottomrule
\end{tabular}
\end{center}

系统有一条设计前提被明确写在接入指南里，值得单独引用，因为它是整套架构所有取舍的源头：

\keybox{\textbf{设计第一原则}\\[0.5em]
任何依赖 Agent 会话中途自觉、或依赖用户主动发起的环节，\textbf{等于不存在}。}

这条前提带来一个可验证的推论：\textbf{必读内容必须放在会话开始的固定动作里}，而不能指望「中途想起来」。整套渐进式加载、启动包、Hook 检查都是为了满足这条约束。

% ======================================================================
\section{分层架构}

系统在纵向上分为\textbf{四层}，职责与可重建性严格区分。

\begin{center}
\begin{tabular}{@{}c@{\hspace{0.8em}}p{0.76\linewidth}@{}}
\toprule
\textbf{层} & \textbf{职责} \\
\midrule
④ 执行 / 治理层 & Hook 策略与宿主适配器；五个固定检查点；硬阻断能力依宿主而异 \\
\addlinespace
③ 工具层 & Python CLI（二十余个子命令）：启动包、检索、保存、交接、提炼、发布、巡检 \\
\addlinespace
② 状态 / 账本层 & \textcolor{WarnRed}{\textbf{不可重建}}：候选队列、操作日志、版本登记、来源快照、费用账本 \\
\addlinespace
① 内容层 & 版本化 Markdown 仓库：权威事实、证据、派生、候选、归档五种材料 \\
\bottomrule
\end{tabular}
\end{center}

\subsection{内容层的五种材料}

内容层按\textbf{可信度与生命周期}分档，而不是按主题分档。

\begin{center}
\begin{tabular}{@{}p{0.16\linewidth}p{0.30\linewidth}p{0.22\linewidth}p{0.22\linewidth}@{}}
\toprule
\textbf{档位} & \textbf{性质} & \textbf{是否权威} & \textbf{新窗口是否默认读取} \\
\midrule
正式事实 & 习惯、知识、雷区、技能、核心热记忆 & \textbf{权威}（须经发布链写入） & 仅压缩投影进启动包 \\
证据 & 会话总结、跨 Agent 日志 & 只证明「当时做过」 & 否，按需追溯 \\
派生 & 记忆包、健康度量、各类索引 & \textbf{可随时重建} & 记忆包必读 \\
候选 & 收件箱、候选队列 & 未经确认，不得当事实 & 否 \\
归档 & 已替代、已撤出材料 & 仅作历史去向 & 否，默认不检索 \\
\bottomrule
\end{tabular}
\end{center}

\textbf{权威与派生副本的边界被显式声明}：系统维护一份「权威源 → 派生副本」的关系清单，并规定工具只报告「发现不一致」，\textbf{不判断谁对}。这是一个刻意的克制——避免工具自动「修正」出一个可能是错的版本。

\subsection{状态层：不可重建的部分}

这是整套系统里最容易被误解、也最需要保护的部分。

\begin{itemize}
  \item \textbf{索引可重建 $\neq$ 账本可丢弃。} 检索索引能从现有 Markdown 重新切分出来，但\textbf{生成不出}候选来源、操作意图、版本历史、来源快照和已消费预算。
  \item 因此候选队列、操作日志、版本登记（写命令的并发控制依据）、来源快照、反馈状态、费用账本被归为\textbf{不可重建}。
  \item 而\textbf{状态根位于仓库之外}，仓库的备份策略\textbf{不自动覆盖它}。系统自己把这一点标记为「当前最大的单点风险」。
\end{itemize}

\subsection{工具层的命令面}

工具层是一个统一的 CLI，命令按「是否写盘」和「是否需要人工授权」分级。

\begin{center}
\begin{tabular}{@{}p{0.19\linewidth}p{0.42\linewidth}p{0.29\linewidth}@{}}
\toprule
\textbf{类别} & \textbf{代表命令} & \textbf{默认行为} \\
\midrule
只读诊断 & 巡检诊断、写入登记对账、宿主接入检查、状态冻结 & 只读，不跑模型、不改数据 \\
读 + 绑定 & 启动包、有界检索 & 只读，输出任务绑定 \\
预览型写 & 索引构建、候选入队、正式发布 & \textbf{默认只预览，须显式确认才真写} \\
阶段保存 & 检查点、交接、单次语义提交 & 携带期望版本号，冲突即拒绝 \\
长时作业 & 日常增量维护、反馈消费、历史提炼 & 受预算与时间窗约束 \\
事实变更 & 正式事实变更发布通道 & 独立通道，可回滚 \\
\bottomrule
\end{tabular}
\end{center}

设计上做了两个防御：

\begin{enumerate}
  \item \textbf{双闸门}——所有写命令默认预览，必须显式加写入开关才真写。这一条从机制上避免了「以为在看，其实在写」。
  \item \textbf{演练模式下不可能产生计费}——干跑时预留一律被拒绝。
\end{enumerate}

% ======================================================================
\section{运转机制：用户实际怎么用}

\subsection{一个会话的完整生命周期}

\keybox{\small
\textbf{一个会话的完整生命周期}\\[0.7em]
\textbf{① 宿主自动读取薄入口}（引导文件，不含规则正文）\\[0.3em]
\textbf{② Hook 检查}入口文件存在性 + 记忆包新鲜度\\
\hspace{1.6em}\emph{只检查，不声称已把正文交给模型}\\[0.3em]
\textbf{③ 执行一次启动引导命令}，返回：约束块 + 任务身份 + 当前目标/下一步 + 适用线索\\
\hspace{1.6em}\textcolor{WarnRed}{\textbf{【不读】}全量规则、首页全文、所有项目状态、完整共享包}\\[0.3em]
\textbf{④ 干活}（需要历史时才发起有界检索，默认最多 5 条）\\[0.3em]
\textbf{⑤ 有效阶段结束时保存一次}（检查点 / 交接）\\
\hspace{1.6em}携带期望版本号 · 核对回执与读回版本\\[0.3em]
\textbf{⑥ 会话结束}：写会话总结；有真实新收获才写候选块}

\subsection{用户实际输入的东西}

用户的日常输入量被压到\textbf{很低}，这是整套设计的目标之一。

\begin{center}
\begin{tabular}{@{}p{0.16\linewidth}p{0.46\linewidth}p{0.24\linewidth}@{}}
\toprule
\textbf{时机} & \textbf{用户 / Agent 动作} & \textbf{频率} \\
\midrule
新窗口 & 一条引导命令（含工作区与宿主名） & 每窗口一次 \\
同窗口再问 & 加「已知规则版本」参数 & 版本未变则零成本 \\
需要历史 & 一条有界检索命令 & 按需 \\
阶段完成 & 一条保存命令 & 每有效阶段一次 \\
换窗口 & 一条交接命令 & 交接时 \\
\bottomrule
\end{tabular}
\end{center}

\textbf{关键点}：用户不需要在每次会话里重复申明偏好、背景和约束——这些由启动包在会话最开始一次性注入。

\subsection{工作区天然交接}

交接不是「把对话记录复制过去」，而是\textbf{把任务身份和现场结构化地移交}。

\begin{itemize}
  \item 每个任务绑定一个三元组：\texttt{（项目标识，任务标识，交接卡路径）}。
  \item 保存时\textbf{携带这套绑定}，因此不会被其他窗口改变焦点后换错卡。
  \item 交接卡带\textbf{修订号}；新窗口读回时必须核对版本。
  \item 冲突时\textbf{拒绝覆盖并给出当前哈希}，要求显式确认接管。
  \item 重复提交同一操作标识 $\rightarrow$ \textbf{幂等去重}（不是丢写），回执明确说明「重复请求已去重」。
  \item 过期版本号写入 $\rightarrow$ \textbf{拒绝}，并报告「当前记录修订号与提交方期望值不符」。
\end{itemize}

这套机制让\textbf{任意 Agent、任意窗口、任意时间}都能接手一个任务，而不依赖「上一个窗口的对话还在」。这正是工作区自然交接的实现基础。

\subsection{焦点与任务的分离}

系统特意区分了两个概念：\textbf{绑定}（这个窗口属于哪个任务，保存时用）与\textbf{焦点}（新窗口默认选哪个任务，\textbf{仅供默认选择}）。

由此产生两条纪律：\textbf{已完成的任务卡不自动重启}；\textbf{只有需要改变新窗口默认任务时才动焦点}。

% ======================================================================
\section{亮点一：多 Agent 协作}

\subsection{共享方式：薄入口，而非复制规则}

多个不同类型的宿主（命令行 Agent、编辑器内 Agent、桌面应用 Agent、独立运行时）接入了同一个仓库。接入方式刻意做成了\textbf{薄入口}：

\begin{itemize}
  \item 每个宿主的全局指令文件\textbf{只做引导，不承载规则正文}。
  \item 正式规则\textbf{只有一份权威}，其余入口全部指向它。
  \item 系统记录了历史教训：早期曾出现「两个宿主各自承载 80--115 行规则正文，与权威规则形成\textbf{三套并存的 schema}」，收敛时花了大力气返工。
\end{itemize}

由此得出一条明确原则：\textbf{薄入口越长，规则漂移越快。}

\subsection{跨 Agent 独立复现的设计意图（注意：此机制未落地，见 §9.1）}

系统在设计层提出了一条很有价值的思路。

\keybox{按\textbf{不同的 Agent 名}判断「有多少独立来源确认过这条经验」。\\[0.4em]
\textbf{单 Agent 重复再多也升不上去}——因为同一个模型说十遍错话仍是错话。}

为了让这个机制有效，接入规范强制要求：Agent 名必须\textbf{唯一、稳定、大小写固定}。如果同一个 Agent 被记成两种写法，系统会把它们当成\textbf{两个独立来源}，从而\textbf{虚增置信度}，让未经真实验证的经验冒充「被多 Agent 确认」。

也就是说，\textbf{命名一致性不是格式洁癖，它直接决定置信度计算的正确性}。

\textbf{但必须指出}：经代码核查，\textbf{「按独立 Agent 来源数升权」在实现里并不存在}——承载它的自动计数组件已随旧采收器一同退役，而设计文档里「反复被独立确认会自动升权」的表述\textbf{没有被同步更新}。这构成了本系统最典型的一类缺陷：\textbf{机制被退役了，而承诺它的文案还留着}（详见 §9.1）。

\subsection{协作的产物是分工，而非全员同工}

系统里沉淀了明确的协作编排模式：由高能力模型负责拆解任务、答疑与审查，由低成本模型在独立工作区施工，只有「理解门、授权、测试、审查」全部通过才允许合并。这使多 Agent 协作从「多开几个窗口」升级为\textbf{有成本模型的分工体系}。

% ======================================================================
\section{亮点二：知识沉淀闭环}

\subsection{五级来源等级}

所有事实性内容必须带来源等级，且\textbf{只有前三级允许进入事实层}。

\begin{center}
\begin{tabular}{@{}l@{\hspace{1.4em}}l@{}}
\toprule
\textbf{等级} & \textbf{含义} \\
\midrule
\texttt{user\_confirmed} & 用户明确确认 \\
\texttt{locally\_verified} & 本地实测验证 \\
\texttt{external\_trusted} & 可信外部来源 \\
\midrule
\multicolumn{2}{@{}l@{}}{\textcolor{WarnRed}{\textbf{—— 以上可进事实层，以下仅作线索 ——}}} \\
\midrule
\texttt{agent\_inference} & Agent 推断（仅线索） \\
\texttt{unknown} & 来源不明（仅线索） \\
\bottomrule
\end{tabular}
\end{center}

配套纪律包括：派生物必须声明来源与时点；未知不能伪装已知；\textbf{历史不能补写假证据或作者}；修改正式事实必须经过发布链。

\subsection{完整沉淀链路}

\keybox{\centering\small
会话中的真实收获 $\rightarrow$ 收获块（收件箱）$\rightarrow$ 候选队列\\[0.35em]
$\downarrow$ 语法校验 + 去重 + 作用域检查\\[0.35em]
机器复核（\textbf{此段需要模型能力}）$\rightarrow$ 暂存发布清单（含审计与回滚记录）\\[0.35em]
$\downarrow$ \textbf{人工策展确认}\\[0.35em]
正式发布 $\rightarrow$ 索引回读（可检索、可引用）}

链路里有几个\textbf{刻意设置的闸门}：

\begin{itemize}
  \item \textbf{语法校验是硬闸门}：格式不合法的收获块\textbf{整块被拒}并进入拒绝区，不会静默丢失。
  \item \textbf{作用域不得降级}：没声明适用范围时，系统\textbf{不会默认放宽为全局}，而是判为「待确认」——宁可当作线索，也不冒充已验证。
  \item \textbf{人工策展不可省}：机器只能写到「暂存发布清单」，正式发布需要人的确认。
  \item \textbf{反馈不自动验证知识}：使用反馈\textbf{只影响匹配场景的召回}，不会自动把一条经验升级为已验证。
\end{itemize}

\subsection{适用范围是一等公民}

检索采用\textbf{多维度作用域模型}：全局 / 特定 Agent / 特定项目 / 特定工具 / 特定场景。

\begin{itemize}
  \item \textbf{不同维度之间是 AND，同一维度多值之间是 OR。}
  \item \textbf{未知范围不放宽为全局。}
  \item \textbf{记录者不等于适用者}——「谁记的」和「对谁适用」是两件必须分开声明的事。
\end{itemize}

这一条解决了很多记忆系统的通病：一条在特定工具下成立的经验，被当成普适真理到处套用。

% ======================================================================
\section{亮点三：雷区避雷}

\subsection{雷区的三档生命周期}

失败经验不是简单入库就生效，而是走一条\textbf{观察期 $\rightarrow$ 生效 $\rightarrow$ 退役}的路径。

\begin{center}
\begin{tabular}{@{}p{0.15\linewidth}p{0.42\linewidth}p{0.31\linewidth}@{}}
\toprule
\textbf{档位} & \textbf{含义} & \textbf{是否注入新会话} \\
\midrule
观察档 & 已入库，证据强度待验证 & 不进提醒，进入观察位轮转 \\
生效档 & 已验证，具备提醒资格 & \textbf{进提醒}（按置信度 $\times$ 严重度排序） \\
退役档 & 已过期或被替代 & 不计入提醒，仅作历史 \\
\bottomrule
\end{tabular}
\end{center}

\subsection{观察位与回报义务}

系统不把「待验证」藏着，而是\textbf{主动暴露并索取反馈}。

\begin{itemize}
  \item 每次会话从观察档里轮转\textbf{固定数量}的条目注入（轮转指针可见）。
  \item 被注入的条目附带\textbf{明确要求}：本次会话如果实际涉及对应工具，\textbf{必须回报}结果。
  \item 回报有四种标准取值：\texttt{avoided}（规避成功）/ \texttt{hit\_anyway}（仍然踩中）/ \texttt{falsified}（被证伪）/ \texttt{irrelevant}（不相关）。
  \item \textbf{命中与证伪都能改进系统}——证伪同样是有价值的信号。
\end{itemize}

这是一个\textbf{闭环}设计：不只是「告诉 Agent 别踩坑」，还回收「这次到底踩没踩」，从而让置信度不是纸面推测。

% ======================================================================
\section{亮点四：上下文 token 耗费低}

这是整套系统在工程上最扎实的部分，有多个\textbf{可核验的数字约束}。

\subsection{启动包预算}

\begin{center}
\begin{tabular}{@{}p{0.30\linewidth}p{0.62\linewidth}@{}}
\toprule
\textbf{项目} & \textbf{值} \\
\midrule
默认预算 & \textbf{1800 token} \\
可选档位 & 精简 1800 / 标准 3200 / 审计 5000 \\
裁剪规则 & 关键块（规则约束、任务卡、当前风险）\textbf{全保留}；可选块按优先级丢弃 \\
超预算行为 & \textbf{显式报告超支，不静默丢弃关键内容} \\
\bottomrule
\end{tabular}
\end{center}

有一个设计决策值得单独记录：评审\textbf{否决了「直接上调默认预算」}，理由是「那会把\emph{启动包设计问题}变成\emph{预算问题}」。同时\textbf{明确不做自动升档}，因为自动升档会让超支\textbf{静默消失}，与「检测失败必须暴露」的原则冲突。

\subsection{分层记忆的体积上限}

\begin{center}
\begin{tabular}{@{}p{0.22\linewidth}p{0.30\linewidth}p{0.40\linewidth}@{}}
\toprule
\textbf{层} & \textbf{上限} & \textbf{注入方式} \\
\midrule
核心热记忆 & $\le$ \textbf{2 KB} & 无条件注入（只放「半年不变级」环境事实） \\
上下文包 & $\le$ \textbf{8 KB}（实测约 6.6 KB） & 会话开始必读，\textbf{单段式注入，不二次读取} \\
检索结果 & \textbf{默认最多 5 条} & 只在需要时发起 \\
\bottomrule
\end{tabular}
\end{center}

核心热记忆有明确的\textbf{内容纪律}：用户偏好一律以习惯库为准，本文件\textbf{不复制}；失败模式进雷区库，\textbf{不重复}。也就是说体积上限不仅靠「删」，更靠\textbf{职责不重叠}。

\subsection{节省 token 的四个机制}

\begin{enumerate}
  \item \textbf{同窗口版本缓存}——规则版本未变则只返回一行版本确认，不重复注入规则正文。
  \item \textbf{渐进式阅读}——导航顺序是从概览到具体，\textbf{信息足够且风险低即停}；只有不足、冲突、过期或高风险时才深入。\textbf{禁止默认全库扫描}（Hook 会对递归整库扫描给出提醒）。
  \item \textbf{有界检索}——最多 5 条，最多一次改写，一次改写仍无答案就说明边界。\textbf{信息够用即停，不等待后台学习。}
  \item \textbf{派生物重建而非持久化}——派生目录整体被版本控制排除，随时可重建，不需要人工维护一致性。
\end{enumerate}

综合效果：一个全新窗口的\textbf{启动上下文成本被压到 1--2 千 token 量级}，同时仍能拿到任务身份、当前目标、适用偏好与风险提醒。

% ======================================================================
\section{亮点五：执行层的强制力}

薄入口负责「告诉 Agent 应该怎么做」，Hook 负责「在关键时点自动检查或阻断」。系统明确声明 Hook \textbf{只是执行层，不得另造一套规则正文}。它只实现\textbf{五项}检查，不是通用安全框架。

\begin{center}
\begin{tabular}{@{}p{0.17\linewidth}p{0.75\linewidth}@{}}
\toprule
\textbf{检查点} & \textbf{行为} \\
\midrule
会话启动 & 检查入口文件与记忆包新鲜度；\textbf{过期或条目为零时提示失效} \\
检索范围 & 发现默认全库递归扫描时给出提示 \\
写入分级 & 正式层写入\textbf{硬拒绝}；只放行候选区、附件区与本 Agent 总结区 \\
总结触发 & 第一次成功状态变化后建立并复用一份总结；\textbf{失败的工具调用不触发} \\
文档超限 & 项目文档超过 20 KB 且继续增长时\textbf{只提醒，不阻止、不回滚} \\
\bottomrule
\end{tabular}
\end{center}

三个工程上成熟的做法：

\begin{enumerate}
  \item \textbf{硬软分级}——硬 Hook（写入分级、凭据读取、禁止整体覆盖）在宿主支持时直接拒绝，且\textbf{自身失败时优先 fail-closed}；软 Hook（扫描提醒、文档超限）\textbf{绝不回滚已经成功的操作}。
  \item \textbf{不谎报能力}——「宿主不支持阻断」或「没有本机可执行宿主做真实烟测」时，只能标记为\textbf{审计 / 兼容模式}，\textbf{不能宣称已形成硬约束}。
  \item \textbf{诚实的能力边界声明}——文档明确写着：正则属于提醒和防误操作，\textbf{不是安全沙箱}；异步任务完成后 Agent 必须自行确认并保存。
\end{enumerate}

% ======================================================================
\section{真实缺陷}

以下全部来自系统\textbf{自身的度量产物与运行日志}，不含推测。按严重度排列。

\subsection{若干被文档承诺、但实现层并不存在的机制（最严重）}

机制已退役或从未实现，而承诺它的文案仍在生效。

\begin{center}
\begin{tabular}{@{}p{0.30\linewidth}p{0.64\linewidth}@{}}
\toprule
\textbf{被承诺的机制} & \textbf{实况} \\
\midrule
按独立 Agent 来源数自动升权 & \textbf{不存在}——计数组件已随旧采收器退役 \\
五级来源等级 & \textbf{基本未实现}，该字段无枚举校验 \\
候选 $\rightarrow$ 已验证路径 & \textbf{无代码路径}：机器复核明文「不授予已验证」 \\
候选队列正式晋升 & \textbf{从未发生}：账本 728 行，晋升事件 \textbf{0 次} \\
付费预算软硬限 & \textbf{已实现但无生产调用点}，护栏是惰性的 \\
\bottomrule
\end{tabular}
\end{center}

\textbf{根源}：自动化沉淀链存在\textbf{两条并行的候选来源}，从未缝合——一条走「收件箱 $\rightarrow$ 候选队列」，另一条走「历史提炼 $\rightarrow$ 暂存区 $\rightarrow$ 机器复核」。工具自己承认：「两个候选源，\textbf{晋升协议断裂}。」

\warnbox{即「知识沉淀闭环」\textbf{账面上是闭环、运行上是断链}。§5 应读作\textbf{设计意图}，不是已跑通的现状。}

\subsection{维护链在报错中运行，提炼积压巨大}

\begin{center}
\begin{tabular}{@{}p{0.38\linewidth}p{0.54\linewidth}@{}}
\toprule
\textbf{指标} & \textbf{值} \\
\midrule
定时维护作业 & \textbf{最后运行结果为失败}（返回码非零） \\
提炼产物 & \textbf{无法解析，判定为不可恢复}（JSON 解析失败） \\
本轮提炼成功数 & \textbf{0} \\
待处理来源队列 & 约 \textbf{1813}（扫描到 1902） \\
单轮实际尝试 & \textbf{2}（配置上限：2 来源 / 240 秒 / 1 块） \\
\bottomrule
\end{tabular}
\end{center}

\textbf{根因}：维护链依赖一个\textbf{本地小参数模型}做结构化提炼，而\textbf{把任意模型输出可靠约束成合法 JSON} 是固有薄弱点。失败模式不是「偶尔慢」，而是\textbf{整轮归零}。按当前吞吐，清空积压需\textbf{约 900 轮以上}——这是\textbf{结构性不匹配}，非偶发延迟。

\subsection{学习闭环事实上很冷}

\begin{center}
\begin{tabular}{@{}p{0.36\linewidth}p{0.30\linewidth}p{0.26\linewidth}@{}}
\toprule
\textbf{指标} & \textbf{实测值} & \textbf{对照} \\
\midrule
雷区条目总数 & 484 & — \\
\textbf{生效档条目} & \textbf{5（占 1.0\%）} & — \\
观察档池 & \textbf{455}（多数「观察 0 次」） & — \\
从未收到任何反馈的条目 & \textbf{98.0\%} & 学术基线 \textbf{44.9\%} \\
反馈中「仍然踩中」/「被证伪」 & \textbf{各 0 条} & — \\
\bottomrule
\end{tabular}
\end{center}

系统的结论原文：\textbf{「观察位加大探索买不到归因」}。直白说：雷区库看似很大，但\textbf{真正能起提醒作用的只有 5 条}；98\% 的条目\textbf{从未产生过一次有效性反馈}——置信度体系\textbf{缺乏真实数据支撑}。

\warnbox{这是「知识沉淀」与「知识可用」之间的真实鸿沟：\textbf{入库量不等于可用量}。}

\subsection{消费链积压，末端必须人工}

\begin{center}
\begin{tabular}{@{}p{0.40\linewidth}p{0.52\linewidth}@{}}
\toprule
\textbf{指标} & \textbf{值} \\
\midrule
待审收获块 & \textbf{66--67}（最老积压 8 天） \\
候选：重复 & 193 \\
候选：因\textbf{作用域未声明}被阻塞 & 4 \\
正式晋升 & \textbf{0} \\
\bottomrule
\end{tabular}
\end{center}

系统自己的免责声明：「\textbf{写入成功 $\neq$ 经验已进入可用记忆}。计数只表示\textbf{材料在场}。」

链路末端（正式晋升）\textbf{必须有模型能力且必须人工策展确认}——无法完全自动化，\textbf{必然成为瓶颈}。

\subsection{数据安全存在双重缺口}

\begin{itemize}
  \item \textbf{版本控制未收口}：未提交变更 \textbf{3000+ 项}，最近一次提交在数日前，\textbf{之后的改动全部只存在于磁盘上}。
  \item \textbf{快照目录为空}：实测 0 文件，且是「从未产生」还是「已被清理」\textbf{未确定}——因此回滚能力\textbf{只是能力判断，不是损失评估}。
\end{itemize}

两者叠加，意味着\textbf{没有可回退的基线}。

\subsection{状态层在仓库外，备份覆盖未核实}

状态根（含\textbf{费用账本}、日志）\textbf{位于仓库之外}，系统自述「\textbf{仓库备份策略不自动覆盖它}」，并把此项列为\textbf{「当前最大的单点风险」}。

\subsection{检索能力的局限}

\begin{itemize}
  \item 索引约 \textbf{1187 文件 / 19122 分节}，检索为\textbf{关键词 $+$ 分节}，\textbf{无语义能力}；「最多 5 条」叠加无向量召回，\textbf{同义改写易漏检}。
  \item \textbf{无真正增量}：索引是\textbf{整库重建}，「增量」实为按时间戳\textbf{触发全量重建}，成本随规模\textbf{线性增长}。
  \item \textbf{两套索引并存}：同用途索引维护了两份（约 102 MiB $+$ 118 MiB）。
  \item 边界声明不一致：「最多 5 条」\textbf{有}代码强制，「最多一次改写」\textbf{没有}；启动包内部会发起 \textbf{5--10 次}检索，且\textbf{静默吞掉}内层告警——只会「少一块内容」，看不到「为什么少」。
\end{itemize}

\subsection{Hook 强制力依宿主而异}

硬阻断\textbf{只在支持相应拦截点的宿主上可靠}；多个宿主的接入被标注为\textbf{「兼容 / 审计模式」}（无本机可执行宿主做真实阻断烟测）。\textbf{「写入治理」目前是分级能力，不是统一保证。}

此外，\textbf{真实多宿主接力从未成功过一次}：静态检查全部通过，而\textbf{真实窗口启动全部失败}。系统表述克制——「零条真实接力，不伪造成功」。\textbf{这是诚实的失败，但也说明该架构「在测试里成立，在真机上未被验证」。}

\subsection{其余落差（择要）}

\begin{itemize}
  \item \textbf{雷区注入实际为 0 条}：484 条中 479 条被过滤，455 条卡在已退役的观察档。
  \item \textbf{坏消息的通知通道是空的}：采收器停摆 8 天，而其告警文件\textbf{全库 0 命中}（从未存在）。
  \item \textbf{同一指标多处对不上}：至少 5 处口径漂移（雷区数相差 74、拒绝数登记 11 现存 7 等）。
  \item \textbf{文档维护负担}：受检 23 个文档、4 个超限，最大 72 KB；策略是\textbf{只提醒、不阻止、不回滚}，故需人工消除。
  \item \textbf{新旧并行}：旧采收器已停用且末次失败；演练残留目录数百个；多份旧版入口以备份形式留存。系统\textbf{明确拒绝「靠删文件来收敛」}，代价是\textbf{读者须自行分辨哪些是当前权威}。
\end{itemize}
\section{缺陷的结构性归因}

收敛为\textbf{四类根因}。

\subsection*{根因一：把确定性职责交给概率性组件}

维护、提炼、机器复核都要求模型产出可解析的结构。失败模式不是「偶尔慢」，而是\textbf{整轮归零}。

\subsection*{根因二：写入产能远大于消费产能}

写入端持续高速产生，消费端（去重、校验、机器复核、\textbf{人工策展}、发布）每步收窄且末端必须人工。结果：\textbf{66 条待审、193 条重复、约 1800 个来源积压、正式晋升 0}。

\subsection*{根因三：观察 / 反馈闭环依赖一个已被自身数据证伪的机制}

瓶颈不在「注入」，而在「回报」。回报依赖会话 Agent 中途自觉——而实测此项执行 \textbf{0 次}。

\subsection*{根因四：没有机制校验「文档承诺的能力是否仍存在」（影响最广）}

系统有很强的文档纪律，却\textbf{不知道自己的哪些部分已经失效}：升权机制退役而文案仍生效、五级等级无实现、日报断更而导航仍列为现行、两套索引并存的说明未收敛。

\textbf{前三条是「某个环节没做到」，这一条是「系统对自身状态失去认知」}——对一个以记忆可信度立身的系统，这是最根本的风险。

% ======================================================================
\section{总结评价}

\subsection{做对了什么}

\begin{enumerate}
  \item \textbf{把「会话起始」作为唯一的可靠注入点}，并据此重建了整个架构。这是从实测推导出的设计，而非从良好愿望推导。
  \item \textbf{权威 / 候选 / 派生 / 证据严格分档}，未知不冒充已知，历史不补写证据，检测失败必须暴露。
  \item \textbf{写入有版本、有登记、有幂等、有冲突拒绝}——交接因此可以真正跨窗口、跨 Agent 可靠传递。
  \item \textbf{token 成本有多重硬约束}（启动包 1800、热记忆 2 KB、记忆包 8 KB、检索 5 条），且采用「显式报告超支」而非「静默偷工」。
  \item \textbf{对自身能力边界异常诚实}——不谎报硬约束、不宣称已恢复完整状态、明确「不是安全沙箱」、真实接力失败时不伪造成功。
  \item \textbf{设计意图层面提出了跨 Agent 独立验证、观察期分级、作用域多维模型等一系列有价值的机制构想}——即使部分尚未落地，其思路本身仍有参考价值。
\end{enumerate}

\subsection{真实的短板}

\begin{enumerate}
  \item \textbf{多条被文档承诺的核心机制在实现层并不存在}，而承诺它们的文案仍在生效——\textbf{它侵蚀的是系统对自身状态的认知能力}。
  \item \textbf{自动化沉淀链押在本地模型的结构化输出上}，失败即整轮归零；且\textbf{存在两条从未缝合的候选来源}，正式晋升计数始终为 0。
  \item \textbf{写入产能远大于消费产能}，积压是结构性的而非偶发的。
  \item \textbf{观察 / 反馈闭环自我矛盾}：依赖一个已被自身数据证伪的机制，导致 98\% 冷反馈、1\% 生效率、零强信号。
  \item \textbf{数据安全双重缺口}：状态层在仓库外且备份未核实，叠加版本控制长期未收口、快照目录为空。
  \item \textbf{检索无语义能力、无真正增量}，且边界声明不一致（部分是规则文本、无代码强制；内部告警被吞）。
  \item \textbf{多宿主协作在测试里成立、在真机上未被验证}（真实接力从未成功过一次）。
  \item \textbf{一致性成本被有意转嫁给人工}，设计上诚实，运营上是长期负担。
\end{enumerate}

\subsection{一句话概括}

\keybox{\justifying
这套系统在\textbf{「如何让多个 AI 共享可靠记忆」}这个问题上，工程化程度相当高——分层、契约、版本、幂等、审计、预算、脱敏、能力边界声明一应俱全，且\textbf{对自身局限的诚实程度远超同类系统}。\\[0.5em]
但它面临两个层次的困难。\textbf{运行层}的困难在闭环末端：\textbf{从「材料入库」到「经验可用」之间，横着一条需要模型能力与人工裁决的窄通道}，目前这条通道的吞吐远低于写入端的产能，且中间还有一条从未缝合的断链。\\[0.5em]
\textbf{认知层}的困难更深：系统\textbf{不知道自己的哪些部分已经失效}。若干部件被退役或从未实现，而承诺它们的文案、规则和导航仍在生效。对一个以「记忆可信度」为立身之本的系统来说，\textbf{这比任何单点技术故障都更根本}。\\[0.5em]
换句话说：\textbf{这套系统已经学会了如何记住，也学会了如何诚实地承认失败；但它还没有学会如何确认「自己声称拥有的能力是否仍然存在」。}}

% ======================================================================
\section*{附录：数据来源}
\addcontentsline{toc}{section}{附录：数据来源}

\begin{center}
\begin{tabular}{@{}p{0.42\linewidth}p{0.50\linewidth}@{}}
\toprule
\textbf{数据类别} & \textbf{来源类型} \\
\midrule
启动包预算、档位与裁剪规则 & 工具层源码中的常量与裁剪分支 \\
分层体积上限、注入规则 & 正式规则文件与派生记忆包 \\
生效 / 观察档计数、冷反馈比、积压量 & 系统自身的健康度量产物 \\
停用、失败、提炼归零、吞吐量 & 定时作业运行日志与状态文件 \\
检索索引规模、并行旧索引 & 巡检结论与运维文档 \\
状态层在仓库外、备份风险、快照为空 & 运维与数据保留文档的自述 \\
Hook 五项检查与硬软分级 & Hook 实现说明与测试用例 \\
中途自觉执行 0 次 & 接入指南中的实测记录 \\
定时作业返回码、端点鉴权状态 & 本机现场实测 \\
\textbf{机制在实现层是否存在、预算护栏有无调用点、晋升事件计数、索引是否增量} & \textbf{工具层源码与账本的逐项代码核查（对文档承诺做交叉验证）} \\
未提交变更数、索引体积、分区一致性 & 现场只读清点 \\
\bottomrule
\end{tabular}
\end{center}

\vspace{0.8em}
\keybox{\small
\textbf{脱敏声明}：本文已剔除全部本机绝对路径、用户名、设备标识、具体项目名称与业务细节。所有比例、计数、体积、预算数字均为系统自身产物的原值，未作估算或修饰。}

\end{document}
